% Options for packages loaded elsewhere
\PassOptionsToPackage{unicode}{hyperref}
\PassOptionsToPackage{hyphens}{url}
\documentclass[
]{article}
\usepackage{xcolor}
\usepackage[margin=24mm]{geometry}
\usepackage{amsmath,amssymb}
\setcounter{secnumdepth}{-\maxdimen} % remove section numbering
\usepackage{iftex}
\ifPDFTeX
  \usepackage[T1]{fontenc}
  \usepackage[utf8]{inputenc}
  \usepackage{textcomp} % provide euro and other symbols
\else % if luatex or xetex
  \usepackage{unicode-math} % this also loads fontspec
  \defaultfontfeatures{Scale=MatchLowercase}
  \defaultfontfeatures[\rmfamily]{Ligatures=TeX,Scale=1}
\fi
\usepackage{lmodern}
\ifPDFTeX\else
  % xetex/luatex font selection
\fi
% Use upquote if available, for straight quotes in verbatim environments
\IfFileExists{upquote.sty}{\usepackage{upquote}}{}
\IfFileExists{microtype.sty}{% use microtype if available
  \usepackage[]{microtype}
  \UseMicrotypeSet[protrusion]{basicmath} % disable protrusion for tt fonts
}{}
\makeatletter
\@ifundefined{KOMAClassName}{% if non-KOMA class
  \IfFileExists{parskip.sty}{%
    \usepackage{parskip}
  }{% else
    \setlength{\parindent}{0pt}
    \setlength{\parskip}{6pt plus 2pt minus 1pt}}
}{% if KOMA class
  \KOMAoptions{parskip=half}}
\makeatother
\setlength{\emergencystretch}{3em} % prevent overfull lines
\providecommand{\tightlist}{%
  \setlength{\itemsep}{0pt}\setlength{\parskip}{0pt}}
\usepackage{mathrsfs}
\usepackage{mathtools}
\DeclareUnicodeCharacter{2212}{\ensuremath{-}}

\DeclareUnicodeCharacter{2265}{\ensuremath{\ge}}
\DeclareUnicodeCharacter{2264}{\ensuremath{\le}}
\DeclareUnicodeCharacter{2295}{\ensuremath{\oplus}}
\DeclareUnicodeCharacter{2297}{\ensuremath{\otimes}}
\DeclareUnicodeCharacter{2208}{\ensuremath{\in}}
\DeclareUnicodeCharacter{2261}{\ensuremath{\equiv}}
\DeclareUnicodeCharacter{21A6}{\ensuremath{\mapsto}}
\DeclareUnicodeCharacter{2192}{\ensuremath{\to}}
\DeclareUnicodeCharacter{03BA}{\ensuremath{\kappa}}
\DeclareUnicodeCharacter{039B}{\ensuremath{\Lambda}}
\DeclareUnicodeCharacter{03C4}{\ensuremath{\tau}}
\DeclareUnicodeCharacter{03A9}{\ensuremath{\Omega}}
\usepackage{bookmark}
\IfFileExists{xurl.sty}{\usepackage{xurl}}{} % add URL line breaks if available
\urlstyle{same}
\hypersetup{
  hidelinks,
  pdfcreator={LaTeX via pandoc}}

\author{}
\date{}

\begin{document}

{
\setcounter{tocdepth}{3}
\tableofcontents
}
\section{Mixed support, joined fibres, and their additive
defects}\label{mixed-support-joined-fibres-and-their-additive-defects}

The support of a vector and its arithmetic value are separate data in
the split-zero programme. This note preserves both, including a nonzero
fibre over the bottom support, noninjective transition maps, and maps
that merge support labels. The mixed terms produced by adding vectors
from different fibres give exact maps into the additive-relation groups.
Independent and folded joins give two explicitly related calculations of
these groups.

The programme source is
\href{https://github.com/KokunoYumeto/zeta-function-research-reader/blob/39ced92952a9c807d51500783b221beb141da3b6/workbenches/splitzero-tandem/tex/support_diagrams.tex\#L9}{\emph{Reconstruction
with changes of support index}, D1--D5},
\href{https://github.com/KokunoYumeto/zeta-function-research-reader/blob/39ced92952a9c807d51500783b221beb141da3b6/workbenches/splitzero-tandem/tex/support_diagrams.tex\#L125}{the
quotient and cohomology comparisons, D6--D7}, and
\href{https://github.com/KokunoYumeto/zeta-function-research-reader/blob/39ced92952a9c807d51500783b221beb141da3b6/workbenches/splitzero-tandem/tex/support_diagrams.tex\#L178}{the
categorical and all-support kernels, D8}. The complete source through
the D8 proof was read for this calculation. Every assertion needed below
is proved for the stated objects.

\subsection{1. The exact support model}\label{the-exact-support-model}

Put \(S=G(\mathbb Z)=\{\tau\}\sqcup\{n^\bullet:n\in\mathbb Z\}\), with
\(e=0^\bullet\). Supported elements add and multiply by integer
arithmetic; \(\tau\) is an external additive identity and multiplicative
absorber. In particular \(\tau\ne e\).

Let \(L\) be a join-semilattice with bottom \(\bot\). For every
\(l\in L\), let \(V_l\) be an abelian group. For \(l\le k\), let
\(\rho_{lk}:V_l\to V_k\) be a group homomorphism, with
\(\rho_{ll}=\mathrm{id}\) and \(\rho_{kh}\rho_{lk}=\rho_{lh}\). Define
\[
M=\mathcal T(L,V)=\coprod_{l\in L}\{l\}\times V_l,
\qquad 0_M=(\bot,0_{V_\bot}),
\tag{MS1}
\] \[
(l,x)+(k,y)=(j,\rho_{lj}x+\rho_{kj}y),\quad j=l\vee k,
\qquad n^\bullet(l,x)=(l,nx),\quad \tau(l,x)=0_M.
\tag{MS2}
\] These formulas define an \(S\)-semimodule. Threefold addition in
either order has support \(l\vee k\vee h\) and amplitude equal to the
sum of the three vectors transported to that fibre. This proves
associativity; commutativity follows from the join and group laws. The
bottom zero acts as the additive identity because
\(\rho_{\bot l}(0)=0\). Supported scalar laws hold in each fibre, and
distributivity over (MS2) follows from the linearity of the transitions.
The scalar \(\tau\) gives the constant global zero map; its
multiplication and distributivity laws follow from that formula and the
additive identity law. These verifications cover every scalar and vector
case.

The scalar \(e\) acts by \[
e(l,x)=(l,0_{V_l}).
\tag{MS3}
\] Thus a nonbottom supported zero \((l,0)\), a nonzero amplitude
\((l,x)\), and the global zero \((\bot,0)\) remain different elements. A
vector \((\bot,x)\) with \(x\ne0\) is also different from the global
zero. No assumption that \(V_\bot=0\) is made.

A morphism of diagrams is a bottom- and join-preserving map
\(\alpha:L\to L'\), with homomorphisms \(f_l:V_l\to W_{\alpha(l)}\)
satisfying \[
f_k\rho_{lk}=\rho'_{\alpha(l),\alpha(k)}f_l\quad(l\le k).
\tag{MS4}
\] It gives the map \[
F_{\alpha,f}:\mathcal T(L,V)\longrightarrow\mathcal T(L',W),
\qquad (l,x)\longmapsto(\alpha(l),f_lx).
\tag{MS5}
\] To check its addition law, both sides have label
\(\alpha(l\vee k)=\alpha(l)\vee\alpha(k)\), and (MS4) identifies their
two transported amplitudes. Its scalar law follows from linearity and
preservation of the bottom zero. Identity and composition are given by
the corresponding pointwise formulas.

For completeness, the exact distinction between two kernels is \[
F_{\alpha,f}^{-1}(0_N)
=\coprod_{l\in\alpha^{-1}(\bot')}\{l\}\times\ker f_l,
\qquad
F_{\alpha,f}^{-1}(eN)
=\coprod_{l\in L}\{l\}\times\ker f_l.
\tag{MS6}
\] The first condition requires both \(\alpha(l)=\bot'\) and \(f_lx=0\);
the second requires only \(f_lx=0\), proving the formulas. The indicated
subsets are subsemimodules: their labels are closed under joins, and
(MS4) takes fibre kernels into fibre kernels. Their inclusions have the
respective kernel and inverse-image universal properties because any map
with the indicated image factors uniquely through the subset. The
inclusion from the first displayed object to the second retains the
actual relationship between them.

\subsection{2. The free module retains the supported
zeros}\label{the-free-module-retains-the-supported-zeros}

Let \[
\mathscr L(M)=\mathbb Z[M]/\mathbb Z[0_M],
\tag{MS7}
\] where \(\mathbb Z[M]\) here means the free abelian group on the
underlying set, and the quotient is a quotient of abelian groups. Write
\(b_{l,x}\) for the class of \([(l,x)]\), so \(b_{\bot,0}=0\). This is
not a quotient ring obtained from the addition of \(M\); no such ring
operation is being assumed.

The coefficient ring \[
A=\mathbb Z[(\mathbb Z,\cdot)]
\tag{MS8}
\] is the monoid ring of the multiplicative monoid of integers,
including its zero basis element \([0]\). It acts by
\([n]b_{l,x}=b_{l,nx}\). Associativity follows from \(m(nx)=(mn)x\), the
element \([1]\) acts as the identity, and the relation \(b_{\bot,0}=0\)
is preserved. The scalar \(\tau\) has zero action and is not the basis
element \([0]\).

Put \(E=[0]\), \(F=1-E\), and \[
s_l=b_{l,0},\qquad r_{l,x}=b_{l,x}-s_l.
\tag{MS9}
\] Then \(s_\bot=0\), \(r_{l,0}=0\), \(Es_l=s_l\), and \(Er_{l,x}=0\).
The following decomposition preserves the exact support index: \[
\mathscr L(M)
=\left(\bigoplus_{l\ne\bot}\mathbb Zs_l\right)
\oplus
\left(\bigoplus_{l\in L}U(V_l)\right),
\qquad
E\mathscr L(M)=\bigoplus_{l\ne\bot}\mathbb Zs_l.
\tag{MS10}
\] Here \(U(V_l)\) is the free abelian group with basis \(r_{l,x}\) for
\(x\in V_l\setminus\{0\}\), and identifies with \(F\mathscr L(M)\)'s
indicated summand. Equivalently \(U(V)\) is the augmentation kernel in
the free group \(\mathbb Z[V]\), with \(r_x=[x]-[0]\).

\textbf{Proof of (MS10).} At each nonbottom label, replace the basis
\(b_{l,x}\) by \(s_l=b_{l,0}\) and \(r_{l,x}=b_{l,x}-b_{l,0}\) for
\(x\ne0\). The inverse basis change is \(b_{l,0}=s_l\),
\(b_{l,x}=s_l+r_{l,x}\). At bottom the surviving basis is already
\(b_{\bot,x}=r_{\bot,x}\) for \(x\ne0\). These changes give bases of
disjoint free summands, proving the direct sum. The formulas for \(E\)
and \(F\) follow on this basis. In particular the entire nonzero bottom
fibre is present in the second summand. \(\square\)

The \(A\)-action on the second summand factors through \[
\Lambda=A/(E)=\mathbb Z[\mathbb Z\setminus\{0\}],
\qquad [n]r_{l,x}=r_{l,nx}\quad(n\ne0).
\tag{MS11}
\] The quotient ring statement follows by removing the basis element
\([0]\): it spans an ideal since \([n][0]=[0]\), and products of nonzero
integers are nonzero. On the first summand every basis scalar \([n]\),
including \([0]\), acts as the identity.

The map induced by (MS5) has exact component formulas \[
\mathscr L(F_{\alpha,f})(s_l)=s'_{\alpha(l)},
\qquad
\mathscr L(F_{\alpha,f})(r_{l,x})=r'_{\alpha(l),f_lx}.
\tag{MS12}
\] These follow by subtracting the images of \(b_{l,x}\) and
\(b_{l,0}\). They preserve the \(E,F\) decomposition and commute with
all scalar actions. If \(\alpha(l)=\bot'\), the first image is zero, but
the second image is nonzero whenever \(f_lx\ne0\). Therefore collapsing
a support label does not automatically destroy its nonzero amplitude.

There is a precise free basis for the kernel of the second map in
(MS12), even when many labels merge. Let \[
P=\{(l,x):l\in L,\ x\ne0\},\qquad
P_0=\{(l,x)\in P:f_lx=0\}.
\tag{MS13}
\] For every attained target pair \(q=(a,y)\) with \(y\ne0\), choose one
\(p_q=(l_q,x_q)\) such that \((\alpha(l_q),f_{l_q}x_q)=q\). Then the
kernel has basis \[
\{r_{l,x}:(l,x)\in P_0\}
\ \sqcup\
\{r_{l,x}-r_{l_q,x_q}:
 (\alpha(l),f_lx)=q,\ (l,x)\ne p_q\}.
\tag{MS14}
\] Its image is the free subgroup generated by the attained target
pairs. To prove this, partition the source basis into the lost set
\(P_0\) and the nonempty inverse-image sets of each attained pair. On
such an inverse-image set the map is the augmentation to the one target
basis vector. Its kernel has the differences from the chosen
representative as a basis, by the same invertible basis change used in
(MS10). The lost basis vectors map to zero independently. The direct sum
over all these disjoint sets proves the kernel and image statements.
Only finite sums occur, so infinite index sets cause no additional step.
For the \(E\) summand, the identical proof gives lost generators \(s_l\)
with \(\alpha(l)=\bot'\) and differences \(s_l-s_{l_a}\) among nonbottom
labels with the same nonbottom image.

\subsection{3. The exact mixed term at a
join}\label{the-exact-mixed-term-at-a-join}

For an abelian group \(V\), define the surjection \[
\epsilon_V:U(V)\to V,\qquad \epsilon_V(r_x)=x,
\qquad K(V)=\ker\epsilon_V.
\tag{MS15}
\] Its kernel is generated by \[
a_{x,y}=r_{x+y}-r_x-r_y\qquad(x,y\in V).
\tag{MS16}
\] Indeed these elements evaluate to zero. In the quotient by the
subgroup they generate, the function \(x\mapsto\overline{r_x}\)
preserves addition, zero, and inverses. It is inverse to the induced
evaluation map because the two composites fix every \(x\in V\) and every
generating class \(\overline{r_x}\). This proves the kernel assertion.
The formula \([n]a_{x,y}=a_{nx,ny}\) also proves that \(K(V)\) is a
\(\Lambda\)-submodule.

For \(l,k\in L\) and \(j=l\vee k\), define \[
\begin{split}
\beta_{l,k}:U(V_l)\otimes_{\mathbb Z}U(V_k)&\longrightarrow K(V_j),\\
r_{l,x}\otimes r_{k,y}&\longmapsto
r_{j,\rho_{lj}x+\rho_{kj}y}
-r_{j,\rho_{lj}x}-r_{j,\rho_{kj}y}.
\end{split}
\tag{MS17}
\] The values are assigned on a tensor product of free bases and
extended bilinearly. A zero vector gives zero on the right, so the
notation also works for \(x=0\) or \(y=0\). Evaluation of the right side
is zero, giving the displayed codomain. This is bilinear in the free
generators \(r_x,r_y\); no assertion that \(x\mapsto r_x\) is additive
is made.

Give the tensor product the diagonal \(\Lambda\)-action
\([n](r_x\otimes r_y)=r_{nx}\otimes r_{ny}\). This is an action because
multiplication of basis scalars composes the two integer
multiplications, and it extends by the monoid-ring law. Equation (MS17)
and linearity of the transitions prove that \(\beta_{l,k}\) is
\(\Lambda\)-linear for this action.

The actual addition in \(M\), embedded into its contracted free group,
is recovered by the exact equation \[
b_{j,\rho_{lj}x+\rho_{kj}y}
=s_j+U(\rho_{lj})r_{l,x}+U(\rho_{kj})r_{k,y}
+\beta_{l,k}(r_{l,x}\otimes r_{k,y}).
\tag{MS18}
\] All terms after \(s_j\) lie in the \(j\)-summand of (MS10); here
\(U(\rho_{lj})r_{l,x}=r_{j,\rho_{lj}x}\). Substitution from (MS17)
cancels the two transported terms and leaves
\(s_j+r_{j,\rho_{lj}x+\rho_{kj}y}\), which is the left side by (MS9).
Thus the mixed term records the exact missing part of addition after the
two amplitudes have been transported to their common support.

\textbf{Naturality, including support collapse.} For a diagram morphism
(MS4), let \(j=l\vee k\). The following equality has source
\(U(V_l)\otimes U(V_k)\) and target \(K(W_{\alpha(j)})\): \[
K(f_j)\,\beta_{l,k}
=\beta'_{\alpha(l),\alpha(k)}
\bigl(U(f_l)\otimes U(f_k)\bigr).
\tag{MS19}
\] Here \(K(f_j)\) is the restriction of \(U(f_j)\), since evaluation
commutes with homomorphisms. On \(r_x\otimes r_y\), both sides are the
additive defect of \(f_j\rho_{lj}x=\rho'_{\alpha(l),\alpha(j)}f_lx\) and
\(f_j\rho_{kj}y=\rho'_{\alpha(k),\alpha(j)}f_ky\). The equality
\(\alpha(j)=\alpha(l)\vee\alpha(k)\) ensures the target labels agree.
This proves (MS19) on a free basis and hence everywhere, without
excluding bottom images.

\textbf{Symmetry and the associativity identity.} Exchanging the two
tensor factors takes \(\beta_{l,k}\) to \(\beta_{k,l}\), since their
amplitudes add commutatively in \(V_{l\vee k}\). For three labels
\(l,k,m\), write \(j=l\vee k\), \(q=k\vee m\), \(h=l\vee k\vee m\), and
define \[
x_j=\rho_{lj}x,\quad y_j=\rho_{kj}y,\qquad
y_q=\rho_{kq}y,\quad z_q=\rho_{mq}z.
\] The following identity lies in \(K(V_h)\): \[
\begin{split}
K(\rho_{jh})\beta_{l,k}(r_{l,x}\otimes r_{k,y})
&+\beta_{j,m}(r_{j,x_j+y_j}\otimes r_{m,z})\\
=K(\rho_{qh})\beta_{k,m}(r_{k,y}\otimes r_{m,z})
&+\beta_{l,q}(r_{l,x}\otimes r_{q,y_q+z_q}).
\end{split}
\tag{MS20}
\] To prove it, put \(a=\rho_{lh}x\), \(b=\rho_{kh}y\),
\(c=\rho_{mh}z\). The first side becomes
\((r_{a+b}-r_a-r_b)+(r_{a+b+c}-r_{a+b}-r_c)\), while the second is
\((r_{b+c}-r_b-r_c)+(r_{a+b+c}-r_a-r_{b+c})\). Both equal
\(r_{a+b+c}-r_a-r_b-r_c\) in the specified \(h\)-fibre. This proves the
identity with all transports and domains retained.

\textbf{Cancellation gives a nonzero mixed term.} If
\(a=\rho_{lj}x\ne0\) and \(\rho_{kj}y=-a\), then \[
\beta_{l,k}(r_{l,x}\otimes r_{k,y})=-r_{j,a}-r_{j,-a}\ne0.
\tag{MS21}
\] If \(a\ne-a\), the two terms are distinct free basis vectors. If
\(a=-a\), the expression is \(-2r_{j,a}\), which is nonzero in a free
abelian group. The sum of the two original vectors is \((j,0)\). This is
a supported zero when \(j\ne\bot\), and is the global zero only when
\(j=\bot\). In either case the additive defect (MS21) remains nonzero.
In contrast, the actual global-zero input \((\bot,0)\) gives
\(\beta_{\bot,k}(0\otimes r_y)=0\); a nonzero bottom-fibre input need
not give zero.

\textbf{Vanishing transitions have typed kernels.} If \(x\ne0\) and
\(\rho_{lj}x=0\), then \(U(\rho_{lj})r_{l,x}=0\), so \(r_{l,x}\) is an
actual nonzero kernel generator. On the unreduced free group the map
instead sends \(b_{l,x}\) to \(s_j\), which is nonzero when
\(j\ne\bot\). Its support map sends \(s_l\) to \(s_j\). The mixed term
with this transported amplitude is zero because (MS17) becomes
\(r_b-r_0-r_b=0\). These are three exact maps with different codomains
and retained data, not an identification of the original vector with
absence.

\subsection{4. Independent joins: the mixed term is a full direct
summand}\label{independent-joins-the-mixed-term-is-a-full-direct-summand}

For any abelian groups \(V,W\), take the diamond of labels
\(\{\bot,l,k,j\}\), with \(l\vee k=j\), and diagram \[
V_\bot=0,\quad V_l=V,\quad V_k=W,\quad V_j=V\oplus W,
\qquad \rho_{lj}(x)=(x,0),\quad \rho_{kj}(y)=(0,y).
\tag{MS22}
\] The other transitions are identities or the unique maps from the zero
group. They satisfy all composition identities. Define \[
c_{x,y}=r_{(x,y)}-r_{(x,0)}-r_{(0,y)},
\qquad x\ne0,\ y\ne0,
\tag{MS23}
\] and let \(X(V,W)\) be the subgroup they generate. There is an exact
direct decomposition \[
U(V\oplus W)=U(V)\oplus U(W)\oplus X(V,W),
\qquad X(V,W)\cong U(V)\otimes U(W),
\tag{MS24}
\] where the first two inclusions are induced by the coordinate
inclusions and the last isomorphism is precisely \(\beta_{l,k}\).

\textbf{Proof.} A basis of \(U(V\oplus W)\) consists of the nonzero axis
generators \(r_{(x,0)},r_{(0,y)}\) and the generators \(r_{(x,y)}\) with
both coordinates nonzero. Replace each latter generator by \(c_{x,y}\).
The inverse change is \(r_{(x,y)}=c_{x,y}+r_{(x,0)}+r_{(0,y)}\). This
gives a basis and proves the direct sum. The tensor product of the free
bases of \(U(V)\) and \(U(W)\) has basis \(r_x\otimes r_y\), and (MS17)
sends it exactly to the basis \(c_{x,y}\) of the third summand. Hence
\(\beta_{l,k}\) is injective and onto that summand. \(\square\)

The projection map \((U(\pi_V),U(\pi_W))\) sends the two axis summands
identically to \(U(V),U(W)\), and kills every \(c_{x,y}\). Therefore \[
X(V,W)=\ker\bigl(U(V\oplus W)\longrightarrow U(V)\oplus U(W)\bigr),
\tag{MS25}
\] and the decomposition also gives \[
K(V\oplus W)=K(V)\oplus K(W)\oplus X(V,W).
\tag{MS26}
\] For (MS26), evaluation is ordinary evaluation on the two axis
summands and zero on the mixed summand. Its first two values lie in
independent coordinates of \(V\oplus W\), so their vanishing is
equivalent to membership in the two individual kernels. This proves the
statement and all inclusions.

In particular, for \(V=W=\mathbb Z\), the elements \(c_{m,n}\),
\(m,n\in\mathbb Z\setminus\{0\}\), form a free basis of the mixed group.
No magnitude or sign has been discarded.

\subsection{5. Folding the join: an explicit nonzero kernel and its
image}\label{folding-the-join-an-explicit-nonzero-kernel-and-its-image}

Replace only the top fibre of (MS22), for \(V=W=\mathbb Z\), by
\(V_j=\mathbb Z\), and replace both incoming transitions by the
identity. The map from the independent diagram to this folded diagram is
the identity on labels, the identity on the two incoming fibres, and \[
\sigma:\mathbb Z\oplus\mathbb Z\to\mathbb Z,
\qquad \sigma(m,n)=m+n
\tag{MS27}
\] on the top fibre. Its two transition squares commute because
\(\sigma(m,0)=m\) and \(\sigma(0,n)=n\). By (MS19), the mixed map for
the folded diagram is exactly \[
U(\mathbb Z)\otimes U(\mathbb Z)
\xrightarrow[\sim]{\ \beta_{\mathrm{ind}}\ }X(\mathbb Z,\mathbb Z)
\xrightarrow{\ U(\sigma)\ }K(\mathbb Z),
\quad
r_m\otimes r_n\longmapsto r_{m+n}-r_m-r_n.
\tag{MS28}
\]

Here is its full kernel and image calculation. Write \[
D=\mathbb Z[z,z^{-1}],\quad I_z=(z-1)D,
\qquad R=\mathbb Z[t,t^{-1},u,u^{-1}],\quad J=(t-1)(u-1).
\tag{MS29}
\] The identification \(U(\mathbb Z)\cong I_z\) sends \(r_n\) to
\(z^n-1\). To prove it, regard \(D\) as the free abelian group on the
basis \(z^n\). Its augmentation kernel has free basis \(z^n-1\),
\(n\ne0\), by the coefficient-sum basis calculation. This kernel equals
\((z-1)D\): for \(n>0\), \(z^n-1=(z-1)(1+\cdots+z^{n-1})\); for \(n<0\),
the same divisibility follows from \(z^n-1=-z^n(z^{-n}-1)\). These
formulas prove both inclusions.

Under this identification, \[
K(\mathbb Z)=I_z^2=(z-1)^2D.
\tag{MS30}
\] Indeed \(\epsilon(r_n)=n\) is the derivative-at-one map on \(I_z\),
since the derivative of \(z^n-1\) at \(1\) equals \(n\), including
negative integers. Every element of \(I_z\) has unique expression
\((z-1)q(z)\), because \(D\) is an integral domain. Its derivative at
one is \(q(1)\). The proved augmentation-kernel calculation says
\(q(1)=0\) exactly when \(q\in(z-1)D\), proving (MS30). The domain
assertion follows, for example, by multiplying two nonzero Laurent
polynomials by powers of \(z\) to make nonzero ordinary polynomials;
their leading coefficients have nonzero product in \(\mathbb Z\).

The tensor product \(I_t\otimes_{\mathbb Z}I_u\) identifies with the
ideal \(JR\), by sending \((t^m-1)\otimes(u^n-1)\) to
\((t^m-1)(u^n-1)\). This is an isomorphism: multiplication by \(t-1\)
identifies the free group \(\mathbb Z[t^{\pm1}]\) with \(I_t\),
multiplication by \(u-1\) does the same for \(I_u\), and the tensor
product of the Laurent monomial bases identifies their tensor product
with \(R\). Under these identifications the stated map is multiplication
by the nonzero element \(J\), which is injective and has image \(JR\).
The same leading-coefficient argument, twice, proves that \(R\) is an
integral domain.

Let \(\delta:R\to D\) be the ring homomorphism
\(t\mapsto z,u\mapsto z\). Then \[
\delta\bigl((t^m-1)(u^n-1)\bigr)
=z^{m+n}-z^m-z^n+1,
\] which is exactly (MS28). Moreover \[
\ker\delta=(t-u)R,
\tag{MS31}
\] because the quotient by \(t-u\) identifies the two invertible
generators and is isomorphic to \(D\), with mutually inverse maps given
by \(t,u\mapsto z\) and \(z\mapsto\overline t=\overline u\). Therefore
the complete comparison is the exact sequence \[
0\longrightarrow J(t-u)R
\longrightarrow JR
\xrightarrow{\ \delta\ }(z-1)^2D
\longrightarrow0.
\tag{MS32}
\] For its kernel, \(\delta(Jq)=(z-1)^2\delta(q)\) vanishes exactly when
\(\delta(q)=0\), since \(D\) is a domain. Equation (MS31) then gives the
left term. For surjectivity, \(\delta(Jq)=(z-1)^2\delta(q)\) runs
through the target because \(\delta\) is onto. The kernel is a free
rank-one \(R\)-module with generator \(J(t-u)\), and has free abelian
basis \(J(t-u)t^a u^b\), \(a,b\in\mathbb Z\), since multiplication by
this nonzero Laurent polynomial is injective. The image has free abelian
basis \((z-1)^2z^c\), \(c\in\mathbb Z\), by the same argument.

For example, \(J(t-u)\) corresponds to \[
r_2\otimes r_1-r_1\otimes r_2,
\tag{MS33}
\] because \(t(t-1)=(t^2-1)-(t-1)\) and \(u(u-1)=(u^2-1)-(u-1)\). Its
folded image is zero by the symmetry of addition, but its
independent-join image \(c_{2,1}-c_{1,2}\) is nonzero by the free basis
in (MS24). Equations (MS27)--(MS33) are the precise map, kernel, and
image connecting these two mixed-support structures.

\subsection{6. Mixed support in complexes and its cohomology
maps}\label{mixed-support-in-complexes-and-its-cohomology-maps}

Let \((C_l^\bullet,d_l)_{l\in L}\) be a diagram of cochain complexes of
abelian groups, with transition cochain maps
\(\rho_{lk}^i:C_l^i\to C_k^i\). At each degree construct
\(M^i=\mathcal T(L,C^i)\), and define \[
\mathfrak d^i:M^i\to M^{i+1},\qquad(l,x)\longmapsto(l,d_l^i x).
\tag{MS34}
\] The transition cochain identities are exactly (MS4), so these maps
are \(S\)-linear. Their composite is the typed supported-zero map \[
\mathfrak d^{i+1}\mathfrak d^i(l,x)=(l,0_{C_l^{i+2}}),
\qquad M^i\longrightarrow M^{i+2}.
\tag{MS35}
\] It is not asserted to be the constant global zero map. On contracted
free groups this is \[
b_{l,x}\longmapsto s_l^{i+2},\qquad
s_l^i\longmapsto s_l^{i+2},\qquad
r_{l,x}^i\longmapsto0.
\tag{MS36}
\] Its source and target are \(\mathscr L(M^i)\) and
\(\mathscr L(M^{i+2})\); the two copies of \(s_l\) are in their stated
degrees.

Consequently the \(F\) summands form the genuine complex \[
\mathscr U_L(C)^i=\bigoplus_{l\in L}U(C_l^i),\qquad
D^i r_{l,x}=r_{l,d_l^i x},\qquad (D^{i+1}D^i)=0.
\tag{MS37}
\] The equality follows on every generator from \(d_l^{i+1}d_l^i x=0\)
and \(r_{l,0}=0\). Its cohomology satisfies \[
H^i(\mathscr U_L(C))=\bigoplus_{l\in L}H^i(U(C_l)).
\tag{MS38}
\] Indeed the differential preserves the direct summands; a
finite-support vector is a cycle exactly when each component is a cycle,
and it is a boundary exactly when each of its finitely many nonzero
components is a boundary. Choosing one preimage for each of those
finitely many components proves the boundary assertion and hence (MS38).

The evaluation maps give the exact sequence of complexes \[
0\longrightarrow\bigoplus_l K(C_l)
\longrightarrow\mathscr U_L(C)
\xrightarrow{\ \bigoplus_l\epsilon\ }\bigoplus_l C_l
\longrightarrow0.
\tag{MS39}
\] Degreewise exactness was proved in (MS15)--(MS16), and evaluation
commutes with each differential because both composites take \(r_x\) to
\(d_lx\). Thus the kernels form subcomplexes and (MS39) is exact as
claimed. Diagram morphisms give the maps prescribed by (MS12), summing
contributions when labels merge. Equations (MS13)--(MS14) remain their
exact degreewise kernel formulas; support merging has not been replaced
by a componentwise injectivity assumption.

There is also a cochain-level mixed map. Define the complex \[
T_{l,k}(C)^i=U(C_l^i)\otimes_{\mathbb Z}U(C_k^i),\qquad
\partial^i=U(d_l^i)\otimes U(d_k^i).
\tag{MS40}
\] This is the tensor product in equal degrees with the displayed
differential; it is not the total tensor product of two cochain
complexes. Its square is zero because
\((U(d_l^{i+1})U(d_l^i))\otimes(U(d_k^{i+1})U(d_k^i))=0\). Applying
(MS19) to the fibre differentials proves \[
\beta_{l,k}^\bullet:T_{l,k}(C)\longrightarrow K(C_{l\vee k})
\tag{MS41}
\] is a cochain map. Explicitly, both sides of the cochain identity take
\(r_x\otimes r_y\) to the additive defect of
\(d\rho_{l,j}x=\rho_{l,j}d_lx\) and \(d\rho_{k,j}y=\rho_{k,j}d_ky\) in
degree \(i+1\) of the \(j\)-fibre. It therefore induces the typed
homomorphism \[
H^i(T_{l,k}(C))\longrightarrow H^i(K(C_{l\vee k}))
\longrightarrow H^i(\mathscr U_L(C)),
\tag{MS42}
\] where the second arrow is induced by the inclusion of that kernel
complex into its labelled summand. Cycles map to cycles and boundaries
to boundaries by the proved cochain identities, which proves that both
arrows are well-defined.

For the independent diagram constructed degreewise from two complexes
\(C,D\), the decompositions (MS24)--(MS26) commute with the
differentials: each axis generator stays on its axis, and
\(c_{x,y}\mapsto c_{d_Cx,d_Dy}\). They consequently give the exact
cochain decomposition \[
U(C\oplus D)
\cong U(C)\oplus U(D)\oplus T(C,D),
\qquad T(C,D)^i=U(C^i)\otimes U(D^i),
\tag{MS43}
\] with the differential of (MS40). On cohomology the third summand is
precisely the image of the independent-join map (MS42); its inclusion is
injective because it has the cochain projection supplied by this direct
decomposition. The same decomposition holds for \(K(C\oplus D)\) with
\(K(C),K(D)\) on its first two summands, by (MS26).

For the folded diagram constructed degreewise from any complex \(C\),
the map \(\beta:T(C,C)\to K(C)\) is onto in every degree, since its
values are exactly the generators (MS16). Its kernel
\(R_\beta(C)^i=\ker\beta^i\) is preserved by the differential by (MS41).
This gives the exact sequence of complexes \[
0\longrightarrow R_\beta(C)
\longrightarrow T(C,C)
\xrightarrow{\ \beta\ }K(C)
\longrightarrow0.
\tag{MS44}
\] Thus the obstruction to injectivity of a mixed join defines an actual
kernel complex with specified maps. For a degree whose group is
\(\mathbb Z\), its groups and maps are exactly the Laurent-polynomial
calculation (MS32), not merely an unspecified kernel. The mixed term,
its support labels, its vanishing transitions, and its folding kernel
are all retained objects of the comparison.

\subsection{7. The global additive quotient and its exact transition
kernel}\label{the-global-additive-quotient-and-its-exact-transition-kernel}

Fibrewise evaluation and imposing the addition of the whole semimodule
have different codomains. Their exact comparison is as follows. Form the
group \[
G_V=\left(\bigoplus_{l\in L}V_l\right)
\Big/
\left\langle\iota_k\rho_{lk}x-\iota_lx:
l\le k,\ x\in V_l\right\rangle,
\tag{MS45}
\] where \(\iota_l\) denotes inclusion into the indicated direct sum
before quotienting. Write \(\eta_l:V_l\to G_V\) for the quotient maps.
They satisfy \(\eta_k\rho_{lk}=\eta_l\). Any family of homomorphisms
\(g_l:V_l\to Q\) with these compatibility identities defines a
homomorphism out of the direct sum that kills all the displayed
relations, and hence factors uniquely through \(G_V\). This proves the
colimit universal property of (MS45) directly.

The canonical map \[
\eta:M\longrightarrow G_V,\qquad (l,x)\longmapsto\eta_l(x)
\tag{MS46}
\] preserves addition because
\(\eta_j(\rho_{lj}x+\rho_{kj}y)=\eta_l(x)+\eta_k(y)\), and sends the
global zero to zero. It is onto: a finite sum of representatives in
(MS45) can be transported to the join of its finitely many labels and
then represented by one vector in that fibre. The empty sum is
represented by the bottom zero.

It is the universal map from the additive monoid \(M\) to an abelian
group. In fact an additive map \(h:M\to Q\) into a group kills every
\((l,0)\), since \((l,0)+(l,0)=(l,0)\) implies \(h(l,0)+h(l,0)=h(l,0)\),
and cancellation in the target group gives \(h(l,0)=0\). Its restriction
\(h_l:V_l\to Q\) is a group homomorphism by the same-fibre instance of
(MS2). For \(l\le k\), the identity \((l,x)+(k,0)=(k,\rho_{lk}x)\) gives
\(h_k\rho_{lk}=h_l\). The colimit universal property then proves the
required unique factorization through (MS46). This argument concerns a
group target and makes no cancellation assumption in \(M\).

Extend (MS46) linearly to \(\overline\eta:\mathscr L(M)\to G_V\). Its
kernel has the exact description \[
\begin{split}
\ker\overline\eta
={}&\bigoplus_{l\ne\bot}\mathbb Zs_l\ \oplus\ Q_V,\\
Q_V
={}&\left(\bigoplus_{l\in L}K(V_l)\right)
+\left\langle r_{k,\rho_{lk}x}-r_{l,x}:
l\le k,\ x\in V_l\right\rangle
\ \,\subseteq\bigoplus_l U(V_l).
\end{split}
\tag{MS47}
\] The first sum on the first line is direct from the second term by
(MS10). The plus sign defining \(Q_V\) does not assert an internal
direct sum. To prove the formula, first quotient the \(F\) piece by
\(\bigoplus_lK(V_l)\), obtaining \(\bigoplus_lV_l\) by (MS15). The
remaining displayed generators then become exactly the relations in
(MS45). The quotient is therefore \(G_V\), and all \(s_l\) already
evaluate to zero. This proves (MS47) without an uncomputed kernel.

The subgroup in (MS47) also equals the subgroup generated by all global
addition relations \[
b_{(l,x)+(k,y)}-b_{l,x}-b_{k,y}.
\tag{MS48}
\] Here the first subscript means the element given by (MS2). One
direction follows because (MS46) is additive. Conversely, the relation
with the two inputs \((l,0),(l,0)\) kills \(s_l\); relations within one
fibre give every generator (MS16) after those support generators have
been killed; and the relation for \((l,x),(k,0)\) when \(l\le k\) gives
every transition difference in (MS47). Thus the quotient by (MS48) kills
all generators of (MS47), proving equality. This identifies the global
additive quotient and the exact maps connecting it to the fibrewise
quotient (MS39).

The colimit has a pointwise description retaining noninjective
transitions: \[
\eta_l(x)=\eta_k(y)
\quad\Longleftrightarrow\quad
\text{there is }h\ge l,k\text{ with }\rho_{lh}x=\rho_{kh}y.
\tag{MS49}
\] For a direct proof, define this relation on the disjoint union in
(MS1). Reflexivity and symmetry follow from the definitions. If two
equalities hold at possibly different upper labels, transport both to
their join to prove transitivity. Addition of two equivalence classes is
defined by transport to the join and addition there; replacing
representatives preserves the result after transport to a common upper
label. The zero classes all coincide, and \((l,-x)\) is the inverse of
\((l,x)\). This constructs an abelian group. A compatible family of maps
out of the fibres gives a unique homomorphism from it, by the formula on
representatives, so it has the same proved universal property as (MS45).
The canonical maps between the two groups fix every fibre representative
and are inverse, proving (MS49).

In particular define \[
N_l=\ker\eta_l
=\{x\in V_l:\text{there is }h\ge l\text{ with }\rho_{lh}x=0\}.
\tag{MS50}
\] This is an exact calculation of the kernel, by (MS49) with the zero
representative. It is a subgroup: witnesses for two elements can be
transported to their join, where their sum and negatives also vanish.
The transitions preserve these groups; if \(x\in N_l\) has a witness
\(h\), then for \(k\ge l\), the image \(\rho_{lk}x\) vanishes after
further transport to \(k\vee h\). Consequently \((N_l)\) is a
subdiagram, and \[
\eta^{-1}(0)=\mathcal T(L,N)
\tag{MS51}
\] as a subset and subsemimodule of \(M\). In this formula the group
target is regarded as \(\mathcal T(\{\bot'\},G_V)\); its bottom fibre is
the whole possibly nonzero group. The support map of (MS46) is the
unique map \(L\to\{\bot'\}\), and its fibre maps are the \(\eta_l\).
These maps satisfy (MS4), so this is an exact instance of (MS5)--(MS6).
The kernel includes supported zeros and all amplitudes killed by later
transitions, with their original labels retained in its source.

\subsection{8. Retaining support beside the global
amplitude}\label{retaining-support-beside-the-global-amplitude}

There is also the comparison with the same support set \[
\Theta:\mathcal T(L,V)\longrightarrow\mathcal T(L,\underline{G_V}),
\qquad (l,x)\longmapsto(l,\eta_lx),
\tag{MS52}
\] where the target diagram has the group \(G_V\) in every fibre and
identity transitions. The equality \(\eta_k\rho_{lk}=\eta_l\) proves
(MS4), and hence \(\Theta\) is \(S\)-linear. Its exact equality relation
is \[
\Theta(l,x)=\Theta(k,y)
\quad\Longleftrightarrow\quad
l=k\text{ and }x-y\in N_l.
\tag{MS53}
\] The target retains the label, so equality first requires \(l=k\); its
amplitude equality is then exactly (MS50). Thus the quotient diagram
\(\overline V_l=V_l/N_l\) has a proved embedding \[
\mathcal T(L,V)\xrightarrow{\ q\ }
\mathcal T(L,\overline V)
\xhookrightarrow{\ \overline\Theta\ }
\mathcal T(L,\underline{G_V}),
\qquad
q(l,x)=(l,[x]),\quad\overline\Theta(l,[x])=(l,\eta_lx),
\tag{MS54}
\] whose composite is \(\Theta\). Well-definedness and injectivity of
the second arrow follow from (MS50); (MS4) proves the transition
identities for both arrows. Its image in the \(l\)-fibre is exactly the
subgroup \(\eta_l(V_l)\subseteq G_V\), and these subgroups increase with
\(l\).

Every induced transition \(\overline\rho_{lk}:V_l/N_l\to V_k/N_k\) is
injective. If \(\rho_{lk}x\in N_k\), there is \(h\ge k\) with
\(\rho_{kh}\rho_{lk}x=0\); the composition identity gives
\(\rho_{lh}x=0\), so \(x\in N_l\). This proves injectivity. The original
\(N_l\)'s and their transition maps remain the kernel diagram of the
degreewise exact sequence \[
0\longrightarrow N_l\longrightarrow V_l\longrightarrow V_l/N_l\longrightarrow0.
\tag{MS55}
\] They have zero colimit: every element of each \(N_l\) becomes zero at
an upper label by definition, and (MS49) applies to their
transition-stable subdiagram. Thus (MS54) identifies exactly the
support-labelled data invisible to the global amplitude, without
declaring that data absent in the original diagram.

The morphism \(\Theta\) is injective exactly when every original
transition \(\rho_{lk}\) is injective. In one direction, injective
transitions make every \(N_l\) zero by (MS50), and (MS53) proves
injectivity. Conversely, injectivity of \(\Theta\) makes every \(N_l\)
zero. If \(\rho_{lk}x=0\), then \(x\in N_l\), so \(x=0\). This proves
that every transition is injective.

There is an exact universal property for (MS54). Let
\((\alpha,f):(L,V)\to(L',W)\) be any diagram morphism whose target
transitions are injective. For \(x\in N_l\), choose \(h\ge l\) with
\(\rho_{lh}x=0\). Naturality gives \[
\rho'_{\alpha(l),\alpha(h)}f_lx=f_h\rho_{lh}x=0.
\tag{MS56}
\] Injectivity of this target transition implies \(f_lx=0\). Every
\(f_l\) therefore factors uniquely through \(V_l/N_l\). The induced maps
satisfy (MS4), since their composites with the surjective quotient maps
do; those equalities determine the induced maps on every class. Hence
\((\alpha,f)\) factors uniquely through the first arrow in (MS54), with
the original support map \(\alpha\). This proves the comparison to
diagrams with injective transitions, while (MS55) retains the exact
diagram removed by that particular comparison.

For an explicit noninjective example, take \(L=\{\bot<a\}\),
\(V_\bot=\mathbb Z\), \(V_a=\mathbb Z\), and \(\rho_{\bot a}=0\). Then
\(G_V\cong\mathbb Z\) via the top fibre: (MS45) kills every bottom
representative and imposes no relation on the top group. Equations
(MS50) give \(N_\bot=\mathbb Z\), \(N_a=0\). The original addition
satisfies \[
(\bot,m)+(a,n)=(a,n),
\qquad
\Theta(\bot,m)=(\bot,0),\quad
\Theta(a,n)=(a,n).
\tag{MS57}
\] In \(\mathscr L(M)\), every \(r_{\bot,m}\), \(m\ne0\), is
nevertheless a nonzero independent basis vector. Its transition image in
the upper reduced group is zero, while the unreduced image of
\(b_{\bot,m}\) is the nonzero supported-zero generator \(s_a\). Thus the
kernel diagram, the free reduced module, the supported-zero component
and the global additive quotient give four explicitly connected views of
this mixed-support calculation.

\end{document}
